\begin{lema}[Reordenación de la contribución \(f6\)]\label{lem:f6}
En la malla de lado \(s\), sean \(\mathcal I{}\) la incidencia casilla--línea de la
sección~\ref{sec:rectangulos}, \(\beta_e=\nu_1(e)/s^2\) y
\(N_v=\mathcal I{}^{\mathsf T}\operatorname{diag}(v)\mathcal I{}\). Los productos de vectores de casillas
son puntuales. Con \(K,L_0,\rho,w\) de la sección~\ref{sec:profile-data}, definamos
\begin{align*}
 D&=L_0-K,& V_{ij}&=K_{ij}(W_{ij}-K_{ij}),&
 \Psi&=\rho^{-1}\mathcal I{}^{\mathsf T}(\beta w),\\
 Z&=\Psi+D\rho\Psi,& h&=w-\mathcal I{}Z,& \alpha&=\beta h,\\
 a&=\mathcal I{}^{\mathsf T}\alpha,& b_\beta&=\mathcal I{}^{\mathsf T}\beta,&
 M&=\mathcal I{}D\mathcal I{}^{\mathsf T}.
\end{align*}
Aquí \(W_{ij}\) es el peso del cruce o cero si éste no existe.
Escribamos \(U(e,j)=\sum_{u\in \familias(e)} K_{uj}-3W(e)\) y
\(t(c,e;j)=M_{ce}+\theta_j+U(c,j)+U(e,j)\).
La cuadratura del perfil físico, con el valor genérico en coincidencias accidentales,
da exactamente \(f6_s=-E_1+E_2\), donde
\begin{align*}
 E_1={}&\tfrac12\operatorname{tr}(N_\beta D N_\alpha D)
       +\tfrac12 b_\beta^{\mathsf T}Va-3b_\beta^{\mathsf T}D\mathcal I{}^{\mathsf T}(\alpha W)\\
     &-\tfrac12\sum_{f\in D_{\rm obl}}a_f
         [(D N_\beta D)_{ff}+(Vb_\beta)_f],\\
 E_2={}&\sum_c\beta_c\sum_{j\in \familias(c)}\sum_{e\in j}
     \frac{\eta_j\nu_1(e)}s\,t(c,e;j)
     \sum_{u\in \familias(e)\setminus\{j\}}Z_u .
\end{align*}
En esta fórmula \(D_{\rm obl}\) reúne las dos familias oblicuas.
\end{lema}
\begin{proof}
Fijado \(c\), la respuesta de segundo orden consta de la fuente local \(\Lambda\) y
la fuente de líneas \(\omega\). Al sumar la respuesta contra \(\beta_e w(e)\), cada
\(u\in \familias(e)\) recibe \([\mathcal I{}^{\mathsf T}(\beta w)]_u=(\rho\Psi)_u\).
La parte integral de la respuesta añade \(\rho D\rho\Psi\) por simetría. La masa
total es, pues, \(\rho Z\). La fuente negativa de \(\omega\) integra \(\Lambda\) en
la fibra de \(u\). En la malla \(\rho_u\eta_u=1/s\), por lo que esa suma sustrae
\(\beta_e(\mathcal I{}Z)_e\) del peso de \(\Lambda(e;c)\). Queda la masa \(\alpha_e\).

Para dos marcas genéricas \(c,e\), la media de \(T_e\) es \(M_{ce}\).
Cada cruce directo tiene varianza \(V_{uj}\), pues su indicador al cuadrado es el propio
indicador. Los cruces distintos tienen transversales distintas en posición genérica.
Así \(\mathbb E T_e^2=M_{ce}^2+\sum_{u\in \familias(e),j\in \familias(c)}V_{uj}\).
Para \(f\in \familias(e)\) la misma cuenta da
\(\mathbb E z_f^2=(\sum_{j\in \familias(c)}D_{fj})^2+\sum_{j\in \familias(c)}V_{fj}\).
Al sumar \(\beta_c\alpha_e\mathbb E\Lambda(e;c)\), el cuadrado \(M_{ce}^2\)
da la traza de \(N_\beta D N_\alpha D\) por multiplicación de índices.
La varianza da \(b_\beta^{\mathsf T}Va\), y el momento lineal con \(W(e)\)
da \(b_\beta^{\mathsf T}D\mathcal I{}^{\mathsf T}(\alpha W)\).
Los cuadrados de \(z_f\) dan \((DN_\beta D)_{ff}+(Vb_\beta)_f\), uno por línea
oblicua de \(e\). Los coeficientes de \(\Lambda\) producen \(E_1\).
En \(N_v\) se incluyen las entradas diagonales \(\sum_{e\in u}v_e\), también cuando
el recorrido vuelve por la misma línea.

En la suma de cruces de \(\omega\), fijemos \(e=u\cap j\), con \(j\in \familias(c)\).
El peso directo por la medida de \(u\) es \(\rho_uK_{uj}=\eta_j\nu_1(e)/s\).
Condicionar el cruce a estar disponible marca entonces \(c,e\) en la fibra \(j\).
En la respuesta media \(M_{ce}\) se repone la diagonal \(\theta_j\).
Los tres cruces con \(j\) en cada extremo llevan ahora el peso cierto \(W\) de la
casilla marcada. Cambiar cada antiguo término \(-K\) a \(-W\) añade \(U(c,j)+U(e,j)\).
Queda \(t(c,e;j)\), multiplicado por \(Z_u\) para cada transversal \(u\ne j\) de \(e\).
La fuente de cruces tiene signo negativo en \(\omega\), y la definición de \(f6\)
tiene otro signo negativo. De ahí \(+E_2\) y \(-E_1\). Esta igualdad de sumas
finitas deja el error espacial de la extrapolación en el estado \status{N}.
\end{proof}
